\documentclass{article}
\usepackage[utf8]{inputenc}
\usepackage[T1]{fontenc}
\usepackage{geometry}
\geometry{
    a4paper,
    left=2.5cm,
    right=2.5cm,
    top=2.5cm,
    bottom=2.5cm
}
\usepackage{graphicx}
\usepackage{float}
\usepackage{hyperref}
\usepackage[english]{babel}
\usepackage{amsmath}
\usepackage{amssymb}
\usepackage{array}
\usepackage{booktabs}
\usepackage{caption}
\usepackage{subcaption}
\usepackage{enumitem}
\usepackage{fancyhdr}
\usepackage{titlesec}
\usepackage{url}
\usepackage{textcomp}
\usepackage{xcolor}
\usepackage{listings}
\usepackage{multirow}
\usepackage{longtable}
\usepackage{tabularx}
\usepackage{times}

% Set up page style
\pagestyle{fancy}
\fancyhf{}
\rfoot{\thepage}
\renewcommand{\headrulewidth}{0pt}

% Title formatting
\titleformat{\chapter}[display]{\normalfont\bfseries}{\chaptertitle\thechapter}{0pt}{\Huge}
\titlespacing*{\chapter}{0pt}{50pt}{40pt}

% Document info
\title{Template Document}
\author{Academic Support Agent}
\date{\today}

\begin{document}
\maketitle

\resumen{
Este documento brinda una plantilla para la preparaci n de trabajos originales que desean ser
publicados como Working Paper bajo certificaci n de la Editorial SOPHIC. Se recomienda que el resumen contenga entre 150 y 200 palabras en un solo p rrafo, donde deben
resumirse el contexto, la motivaci n, la metodolog a empleada, los aportes m s originales, los
resultados y las conclusiones de su trabajo. No deben incluirse citas bibliogr ficas y se recomienda no introducir acr nimos ni f rmulas en el
resumen o en el t tulo. No haga referencias a figuras o a tablas. Como recomendaci n general, escriba
su art culo insertando y eliminado texto a partir de este documento. De esta forma, le ser  m s f cil
respetar los estilos predefinidos. No modifique el archivo cls, solamente modifique el archivo main.tex}


\palabrasclave{
Primera palabra o frase clave, segunda palabra o frase clave, tercera palabra o frase clave. 
M ximo 5 palabras clave}

% Insert here the abstract in English language
\abstract{
Write in English the same text inserted in ``Resumen''.}

% Insert here the keywords of your work in English language
\keywords{
Write in English the same text inserted in ``Palabras clave''.}

\maketitle

\printcontactdata

\section{Introduction}

La Sociedad de Doctores e Investigadores de Colombia invita a los interesados en publicar sus documentos de trabajo (\textit{working papers}) con las siguientes especificaciones y haciendo uso de la plantilla suministrada en este documento. 

 Qu  son los working papers?  Los documentos de trabajo son documentos preliminares de car cter t cnico o cient fico. Usualmente los autores elaboran documentos de trabajo para compartir ideas acerca de un tema o para recibir realimentaci n previa a una presentaci n formal con la comunidad cient fica o para publicar en una revista cient fica. Los documentos de trabajo son a menudo la base para otros trabajos relacionados y pueden ser citados por evaluaciones realizadas por pares  (\cite{youdeowei_scientific_2012}).

 C mo se diligencia el producto en el CvLac?
Apropiaci n social del conocimiento y Divulgaci n p blica de la Ciencia - Circulaci n de conocimiento especializado - Documento de trabajo

\textbf{Tip:} En la introducci n busca exponer los aspectos m s relevantes y representativos del documento de trabajo que se presenta, incluyendo el o los objetivos que se desarrollan en el documento

\section{Methodology}

 Cu l es la estructura y elementos constitutivos de cada documento de trabajo?

\begin{itemize}
  \item El documento de trabajo deber  tener una extensi n m nima de 3000 palabras y que no supere las 11000 palabras (sin incluir tablas, figuras y referencias bibliogr ficas y/o bibliograf a).
  \item El resumen y palabras clave deben estar en espa ol y en ingl s. El idioma del t tulo debe corresponder al cuerpo del documento
  \item Si va a enviar su documento en formato \LaTeX\xspace  por favor env e una carpeta comprimida con los archivos de compilaci n. No se recibir n documentos en formato PDF U otros formatos. Tambi n puede compartir su paquete de compilaci n a trav s de la plataforma \textbf{Overleaf} si lo prefiere.
  \item No modifique el archivo de clase \textit{sophicart.cls}.
  \item El documento debe ser in dito. 
  \item Toda informaci n suministrada debe ser de autor a de la persona que env a el documento, en caso tal que haya sido resultado de un proyecto en conjunto, debe indicarlo en la afiliaci n, agradecimientos, cr ditos y autores. 
  \item Debe incluir el c digo ORCID de los autores. Si no lo tiene, debe crearlo en \url{https://orcid.org/}
  \item Esta plantilla ya est  configurada en formato APA. Solamente ingrese sus referencias en formato \textbf{biblatex} en el archivo \textbf{referencias.bib}
  \item A pesar de todos estos detalles y cuantos otros que se podr an dar, se recomienda nuevamente escribir su art culo copiando, pegando y reemplazando texto a partir de este mismo documento. Esta es la forma m s f cil y segura de respetar los estilos definidos. Por favor, no redefina ning n elemento del estilo (tipograf a, espacios entre textos, m rgenes, u otras medidas definidas en el archivo de clase \textit{sophicart.cls}).
  
\end{itemize}

\textbf{Tip:} En la metodolog a se busca indicar aspectos relacionados con el dise o del trabajo, o forma de abordar el trabajo cuando el trabajo no incluye un componente metodol gico espec fico. 


\section{Results}

La estructura general que se espera para su manuscrito abarca secciones como: introducci n,
materiales, m todos, resultados, discusi n, conclusiones, trabajos futuros, agradecimientos y 
referencias. Estos t tulos pueden combinarse de a dos en una misma secci n y los t tulos trabajos 
futuros y agradecimientos son totalmente optativos. Es com n que la secci n de m todos lleve otro 
t tulo m s relacionado con el aporte original del art culo, pero las restantes secciones se presentan
con los t tulos antes listados. Si existieran demostraciones u otros desarrollos matem ticos extensos,
se recomienda agruparlos en material suplementario.

A continuaci n se dar n m s detalles acerca de las secciones del documento y los formatos para 
insertar los distintos tipos de objetos, como ecuaciones, figuras y tablas.

En el formato de esta publicaci n las secciones y subsecciones del documento no se numeran y 
se insertan con los comandos tradicionales de \LaTeX, es decir
\begin{verbatim}
  \section{Nombre secci n}
  \subsection{Nombre subsecci n}
  \subsubsection{Nombre subsubsecci n}
\end{verbatim}
El formato para los p rrafos ya incluye una sangr a autom tica en la primera l nea y un espacio extra para la separaci n entre p rrafos.

\subsection{Equations}
Las ecuaciones menores o definiciones de variables pueden insertarse directamente en la l nea del
p rrafo, por ejemplo, consid rese que se desea definir una historia $\vect{h}_i^n = w_{i-1}. w_{i-2},
\dots, w_{i-n+1}$ asociada a un s mbolo $w_i$. Observe que una manera sencilla de asegurar la 
uniformidad en el estilo de las ecuaciones es escribir las formulaciones matem ticas siempre en el
entorno correspondiente, es decir, utilizando \verb!$a + b$! para escribir por ejemplo $a + b$
(no escribir directamente como el texto a + b). Por otro lado, recuerde que las unidades de medici n
deben escribirse siempre en formato \emph{redonda}, de modo que  stas no se confundan con variables
(por ejemplo, debe escribirse $1\,\text{m} = 100\,\text{cm}$ en lugar de $1m = 100cm$).

Para insertar ecuaciones m s complejas o que deban ser referenciadas, se recomienda utilizar los
entornos de ecuaciones disponibles en el paquete \verb!amsmath!, recordando que la orden 
\verb!\begin{equation}! numera autom ticamente las ecuaciones. Para escribir ecuaciones sin
numeraci n utilice \verb!\begin{equation*}! o simplemente \verb!$$<expresi n>$$! para
obtener la expresi n en una l nea aparte, por ejemplo $$ \frac{a + b + c}{2} = d. $$
En el caso de una ecuaci n numerada, se debe definir su \emph{etiqueta} con el comando
\verb!\label{ec-1}!:
%----- Ecuaci n.
\begin{equation} \label{ec-1}
 P_l(w_i|\vect{h}_i^{k}) = \sum_{j=0}^{k-1} \lambda_j \hat{P}(w_i|\vect{h}_i^{j}).
\end{equation}
%----- Fin. Ecuaci n.
Para hacer referencia a esta ecuaci n desde el texto debe utilizarse el comando \verb!\eqref{ec-1!,
el cual coloca autom ticamente el n mero de la ecuaci n entre par ntesis. Por ejemplo, ``en
ec.~\eqref{ec-1} se puede ver la estimaci n de la probabilidad de una historia a partir de la simple
combinaci n lineal de historias de orden inferior''.
Recuerde que el uso de \emph{entrecomillados} en \LaTeX\xspace debe realizarse utilizando los comandos
correspondientes, es decir mediante \verb!``texto''! para obtener ``texto''.

\subsection{Figures}
Las figuras deben incluirse debidamente referenciadas utilizando los comandos de \LaTeX\xspace tradicionales,
y nunca deben colocarse como elementos sueltos dentro del texto.
El pie o ep grafe de figura se coloca autom ticamente utilizando el entorno
\begin{verbatim}
\begin{figure}[!tb]
  \centering
  \includegraphics[<options>]{<file>}
  \caption{Ep grafe} \label{<etiqueta>}
\end{figure}
\end{verbatim}
y rellenando el campo correspondiente en \verb!\caption{<>}! (ver Fig.~\ref{fig-1}). Las figuras pueden contenerse en archivos PDF, JPG o PNG, entre otros. Dentro del campo \verb![<options>]! se puede utilizar el modificador
\verb![width=.8\columnwidth]! si es necesario ajustar el tama o de la figura, como ejemplo se ha fijado el
factor \verb!.8!. \textbf{ El comit  acad mico se reservar  el derecho de solicitar que sus im genes cumplan un m nimo de resoluci n y calidad de imagen. Usualmente im genes con un m nimo de 300 dpi se consideran con buena resoluci n.}

%----- Figura.
\begin{figure}[!tb] 
 \centering
 \includegraphics[width=.8\columnwidth]{figura1} 
 \caption{Esquema del mapa utilizado indicando las diferentes formas de reinyecci n y el efecto
          del ruido.} \label{fig-1}
\end{figure}
%----- Fin. Figura.

Preferentemente, las figuras deben disponerse al comienzo o al final de una columna de texto (para lo
cual se hace uso de la opci n \texttt{[!tb]}). No se recomienda disponer las figuras al final del trabajo.
No incluya saltos ni espacios adicionales en los extremos de las figuras ya que  stos se encuentran
debidamente definidos en el archivo de clase.
Si en la figura se utilizan ejes cartesianos, recuerde siempre describir a qu  corresponde cada
eje (etiquetas) con una fuente de tama o no menor a 7 pt para facilitar la lectura.
Para hacer referencia a una figura se debe utilizar la forma abreviada Fig. seguida
del comando \verb!\ref{etiqueta}!, salvo cuando est  al comienzo de un p rrafo, caso en que se
deber  utilizar la palabra completa.

%----- Figura.
\begin{figure*}[!tb] 
 \centering
 \includegraphics[width=.8\columnwidth]{figura2.png} 
 \caption{Ejemplo del \emph{aliasing} que se produce en una cuadr cula con $N = 8$ nodos. Ambos modos
          ($k = 2$ y $k = 10$) toman los mismos valores en los puntos de la cuadr cula.} \label{fig-2}
\end{figure*}
%----- Fin Figura.

Tenga en cuenta que los gr ficos vectorizados brindan una mejor calidad electr nica y de
impresi n, por lo tanto, inserte todas las gr ficas con alg n formato vectorizado o bien, si se
tratase de una fotograf a o imagen m s compleja, utilice formatos con compresi n sin p rdida de
informaci n (se pueden configurar los formatos JPG, PNG, TIF,GIF, etc).
Para incluir figuras que requieran ser visualizadas en el ancho total de la p gina debe utilizarse
nuevamente el entorno \verb!\begin{figure*}!, como en el caso de la Fig.~\ref{fig-2}. 

\textbf{Tip: El comit  acad mico se reservar  el derecho de solicitar que sus im genes cumplan un m nimo de resoluci n y calidad de imagen. Usualmente im genes con un m nimo de 300 dpi se consideran con buena resoluci n.}


\subsection{Tables}
Es preferible que las tablas se dise en utilizando los comandos de \LaTeX\xspace definidos para tal efecto,
y no que se inserten como archivos de im genes ya que esto va en detrimento de la calidad del
documento. Sin embargo, esta opci n puede resultar aceptable si la inserci n se hace a partir de un
formato vectorizado que respete el tama o y el estilo de la tipograf a.
El ep grafe de las tablas es marcadamente diferente al pie de las figuras. En este caso se coloca por
arriba de la tabla, con fuente de tama o 8 pt y p rrafo centrado.
Al igual que en las figuras, es preferible que las tablas se encuentren al principio o al pie de una
columna. El tama o del texto dentro de las tablas no deber a ser inferior a 7 pt ni mayor a 10 pt (el
tama o de letra utilizado en el texto general).
Un ejemplo de este estilo puede verse en la Tabla~\ref{tabla-1}.

\begin{table}[!b]
 \centering
  \caption{Resultados finales y reducci n relativa de los errores (promedios sobre 10 particiones de
 entrenamiento y prueba).} \label{tabla-1}
 {\small
 \begin{tabular}{ccccc}
  \hline
  \hline
  \thead{Errores de \\ reconocimiento} & \thead{SER \\ \%} & \thead{WER \\ \%} & \thead{WAER \\ \%} &
                        \thead{Reducci n \\ \% WER} \\
  \hline
  \hline
  Referencia & $38.30$ & $7.54$ & $8.53$ & $-$ \\
  \hline
  HMM-PASS & $30.55$ & $5.36$ & $6.67$ & $28.91$ \\
  1-PASS & $25.50$ & $4.76$ & $5.70$ & $36.87$   \\
  \hline
  \hline
 \end{tabular}}
\end{table}

\textbf{Tip: En esta secci n se busca exponer los hallazgos principales de la investigaci n realizada. Esta secci n puede ir subdividida en secciones y puede ir unida a la discusi n si as  lo prefiere.} 

\section{Discussion}

\textbf{Cesi n de derechos y autorizaci n de publicaci n:} Por medio de ese formato las personas autoras del documento de trabajo permiten que SoPhIC publique y difunda sin  nimo de lucro la publicaci n, sin que ello represente pago alguno por parte de la asociaci n. Descargar formato y pasos para postulaci n en \url{https://congreso.sophicol.org/publicaciones-y-productos/}. 

\textbf{Proceso:}

\begin{itemize}
\item Postulaci n de documento de trabajo. Enviar del trabajo completo (Documento de PDF si utiliz  Word, o paquete de compilaci n si utiliz  \LaTeX\xspace) de acuerdo a los lineamientos. Encargado: Autores. Instrucciones para subir a Zenodo las puede encontrar en \url{https://tinyurl.com/fj7c7fpa}.
\item Revisi n de calidad. Se evaluar : * cumplimiento editorial de forma, * confirmaci n de trabajo in dito, \textbf{el documento no puede superar un  ndice de 15\% similitud (citas textuales) con ninguna otra publicaci n}. Encargado: Comit  Cient fico SoPhIC. 
\item  Notificaci n de aceptaci n o rechazo a autores. Encargado: Equipo Editorial SoPhIC. 
\item  Env o de formato de cesi n de derechos. El formato lo puede descargar aqu : \url{http://u.pc.cd/n4DrtalK}. Por favor enviar el formato completo a \href{mailto:editorial@sophicol.org}{editorial@sophicol.org}. Encargado: Autores.
\item  Confirmaci n de Digital Object Identifier (DOI) en cat logo ZENODO y de la Editorial de SoPhIC. Encargado: Equipo Editorial SoPhIC. 
\end{itemize}

\section{Conclusion}
Exponer los descubrimientos a la luz de los objetivos y la metodolog a o abordaje del tema planteado inicialmente. Tambi n puede incluir los aspectos pendientes por desarrollar en la investigaci n. 

\section{Acknowledgements}
Reconocer, opcionalmente, las contribuciones de personas o instituciones que hayan tenido un impacto directo en el trabajo, incluyendo fuentes de financiaci n. 

\printbibliography %Prints bibliography
\end{document}